% Standalone article fragment. Requires amsmath, amssymb, amsthm, enumitem.
% Conditional on the stated arithmetic inputs, especially Sun, Theorem 4.1.
\section{A definition with forty existential variables}

\newtheorem{sdtheorem}{Theorem}
\newtheorem{sdlemma}[sdtheorem]{Lemma}
\newtheorem{sdproposition}[sdtheorem]{Proposition}
\newcommand{\sdQ}{\mathbb Q}
\newcommand{\sdZ}{\mathbb Z}
\newcommand{\sdTr}{\operatorname{Tr}_{F/\sdQ}}
\newcommand{\sdord}{\operatorname{ord}}
\newcommand{\sdnum}{\operatorname{num}}
\newcommand{\sdSq}{\operatorname{Sq}(\sdQ)}

The following construction uses the height and pole-predicate inputs of
Sun's preprint, together with its simultaneous prime-pattern theorem
\cite[Theorem~4.1]{Sun2026}. The latter is an additional arithmetic input;
it is not a consequence of the statement that $\operatorname{HTP}_{86}(\sdQ)$
is undecidable. The construction defines the actual subset $\sdZ\subseteq\sdQ$.

\subsection{The arithmetic inputs and fixed constants}

Put $F=\sdQ(\sqrt2)$, let $\sigma$ be its nontrivial automorphism, and take
\[
 E:y^2=x^3-d^2x,\qquad d=75-53\sqrt2.
\]
Fix an even positive integer $m_0$ and a nontorsion point $P_0$ such that
$m_0E(F)=\sdZ P_0$. We use the following data.

\begin{enumerate}[label=(\roman*)]
\item Five distinct rational numbers $\theta_0=0,\theta_1,\ldots,\theta_4$,
an integer $c\ge1$, and a fixed exceptional set of rational primes.
For every fixed enlargement $S$ of that set, there is $H\ge1$ such that
\begin{equation}\label{sd:height}
 h(t)\le H M_S(t)^c \qquad
 (t\in\sdQ\setminus\{\theta_0,\ldots,\theta_4\}).
\end{equation}
Here $M_S(t)$ is the product of primes outside $S$ at which one of the
five finite contacts has positive odd normalized depth. More precisely,
for primitive integral coordinates $t=(u:v)$, use primitive integral
forms proportional to $u-\theta_jv$ and subtract
$\min\{v_p(u),v_p(v)\}$. There is \emph{no contact at infinity}.
This is the coordinate change and height bound of
\cite[Section~2.3]{Sun2026}.

\item A finite multiplier list $\mathcal M\subset\sdQ^\times$, a fixed set
$S_0$ containing $2,3$ and the supports of its elements, a number
$\epsilon>0$, and integers $b_p\ge1$ for $p\in S_0$ with the following
property. For $B\ne0$, define $\Phi_0(B)$ by a finite choice
$\mu\in\mathcal M$, two rational structural witnesses $e,z$, and
\begin{align}
 B'&=\mu B,& J_B&=B'(B'-B^2),\\
 BeJ_Bz&\ne0,&
 B'(B'-e)(B'+3e)&\in\sdSq,\\
 &&B'(B'z^2-eJ_B)(B'z^2+3eJ_B)&\in\sdSq.
 \label{sd:phi}
\end{align}
Then $\Phi_0(B)$ excludes negative odd valuations outside $S_0$.
Conversely it holds whenever $B$ has no negative odd valuation at any
prime, $0<B<\epsilon$, and $v_p(B)\ge b_p$ for every $p\in S_0$.
This is \cite[Proposition~2.3]{Sun2026}. It uses the full square-class
realization input of the H10 manuscript; the weaker prime-squareclass
input isolated in our third version is not sufficient for this application.

\item Given distinct integers $a_1,\ldots,a_k$, an odd squarefree integer
$M$ with $7\nmid M$, a denominator dilation, and arbitrary lower bounds,
there are positive integers $u,w$ for which every $u+a_jw^2$ is prime,
the quotient $u/w^2$ is arbitrarily large, and the prescribed dilation
divides $w$. The selected primes split in $F$ and have two reduction orders
$d_+(q),d_-(q)$ satisfying
\begin{equation}\label{sd:orders}
 v_2(d_+(q))=2,\quad v_2(d_-(q))=6,\quad
 \gcd(d_+(q),d_-(q'))=4
\end{equation}
for every pair of selected primes, including an equal pair. Both orders
are coprime to $M$ and to $7$. These are consequences explicitly included
in \cite[Theorem~4.1]{Sun2026}; the last assertion at $7$ follows from
the prescribed Gaussian residue $1+3i\pmod7$, giving orders $6,2\pmod7$.
\end{enumerate}

We also use the usual formal group and canonical height of $E$.
Write
\[
 z(Q)=-x(Q)/y(Q),\quad \rho(Q)=1/x(Q)\quad(Q\ne O),
 \qquad z(O)=\rho(O)=0.
\]
These functions are defined on every cyclic point: a nonzero cyclic point
is not torsion, so $x(Q)y(Q)\ne0$. Write $\ell_K$ for the formal logarithm
truncated at degree $K$. At a good unramified prime $p>K+1$, two cyclic
points $aP,bP$ in the formal kernel, with $b\ne0$, satisfy
\begin{equation}\label{sd:formal}
 v\!\left(\frac{z(aP)}{z(bP)}\right)=v_p(a/b),\qquad
 \frac{\ell_K(z(aP))}{\ell_K(z(bP))}\equiv a/b\pmod{p^K}
 \quad\text{if }v_p(a/b)\ge0.
\end{equation}
The valuation of the truncated ratio is always $v_p(a/b)$.
At arbitrary fixed places the first identity holds on a sufficiently deep
formal subgroup. These are \cite[Lemma~2.1 and Section~8.1]{Sun2026}.

Choose $c_*>0$, divisible by $64$, so that $c_*P_0$ belongs to a deep
formal subgroup at all places over $S_0$, with $v(z(c_*P_0))\ge2$.
Set $K=2c$, choose $\xi\in\sdQ\setminus\{\theta_j\}$, and fix an integer
$d_1>0$ such that
\[
 r(T)=d_1^2\prod_{j=0}^4(T-\theta_j)\in\sdZ[T].
\]
Fix an integer $b_*>\max\{1,\theta_0,\ldots,\theta_4\}$.
Choose $S$ to contain $S_0$, all bad and ramified primes, the primes at most
$K+1$, and the supports of $c_*,d_1$, all contact coefficients and differences,
and all $\theta_j-\xi$. Fix $H$ in \eqref{sd:height} for this $S$.

Let $C_E\ge0$ satisfy
$|h(x(Q))-2\widehat h(Q)|\le C_E$ on $E(\overline{\sdQ})\setminus\{O\}$.
Fix an integer $L\ge\lceil H\rceil+9$ and put
\[
 C_L=2L^2+6L+6,\qquad
 g_L(S,T,W)=T^2-LS(LS-2)TW+(LS-2)^2W^2-1.
\]
Replace $P_0$ by $P=2^aP_0$ for a fixed nonnegative integer $a$, and put
$m=2^am_0$, so that $G=mE(F)=\sdZ P$ and
\begin{equation}\label{sd:C}
 C:=\frac{2H+C_E+\tfrac32\log2}{2\widehat h(P)}
       <\frac1{2\sqrt{C_L}}.
\end{equation}
No exceptional primes need to be added. In particular $c_*P$ remains in
the chosen formal subgroups.

We require one more pair of fixed positive integers $A,T_0$.
Choose $A$ supported on $S_0$ so that
\begin{equation}\label{sd:A}
 2v_p(A)\ge b_p+2v_p(2)+2\qquad(p\in S_0).
\end{equation}
To choose $T_0$, let $\mathcal G$ be the closure of
$\sdZ(P,P^\sigma)$ in $E(\mathbb R)\times E^\sigma(\mathbb R)$.
Its identity component is a torus and projects onto the identity circle
in each factor. Indeed each projection of $\mathcal G$ is an infinite
compact subgroup of a real elliptic curve, and $\mathcal G$ has finitely
many components. In every component choose a nonempty relatively open set
on which, at both embeddings,
\begin{equation}\label{sd:goodsets}
                 |\rho(Q)|,\ |z(Q)|,\ |1/z(Q)|\le C_0
\end{equation}
for one fixed $C_0\ge1$. Such sets exist because the forbidden poles and
zeros lie over finite subsets of either elliptic factor. Put $C_1=C_0^2$,
and choose $T_0$ divisible by every prime of $S_0$ with
\begin{equation}\label{sd:T0}
 T_0^4>A^2(1+C_1^2)^6/\epsilon.
\end{equation}
All quantities fixed so far are coefficients or fixed rational functions,
not quantified variables. Numerical values have not been emitted here.

\subsection{The aggregation and the finite formula}

For $R\in\sdQ^\times$, field expressions $D_j\in F$, and positive integers
$k_j$, let $\mathcal A(R;(D_j,k_j)_j)$ denote
\begin{equation}\label{sd:agg}
 W=\prod_j\bigl(1+R(D_j/R^{k_j})^2\bigr),\qquad
 \Theta=\sdTr(W^{-1}),\qquad W\Theta\ne0,\qquad
 \Phi_0(R/\Theta^2).
\end{equation}
The displayed quantities are computed expressions. Each occurrence has
only the two structural witnesses of $\Phi_0$ and its two finite square-root
choices. If $p\notin S$ and $v_p(R)$ is odd, this condition implies that
$v_p(R)>0$ and, at some place $v\mid p$,
\begin{equation}\label{sd:agg-sound}
                       v(D_j)\ge k_j\quad\text{for every }j.
\end{equation}
For completeness we will verify the relevant valuations directly.
For soundness, put $\alpha=v_p(R)$ and
$\lambda_v=v(W)=\sum_j\min\{0,\alpha+2(v(D_j)-k_j\alpha)\}\le0$.
If $\lambda=\max_{v\mid p}\lambda_v$, then
$v_p(\Theta)\ge-\lambda$. The odd valuation of $R/\Theta^2$ is nonnegative,
so $\alpha>0$ and $\lambda\ge-\alpha/2$. At a place realizing $\lambda$,
each summand is at least $-\alpha/2$. Hence
$v(D_j)\ge(k_j-3/4)\alpha>k_j-1$; integral valuations give
\eqref{sd:agg-sound}. Zero expressions cause no exception.

For an input $X\in\sdQ$, put
\[
 \beta=X^2+b_*,\qquad r=r(\beta),\qquad N=1+T_0^2r^2,
 \qquad R_0=A^2r.
\]
In particular $r,N,R_0$ are positive. Quantify four nonzero points
$Q_n,Q_s,Q_t,Q_w\in G$, a point $H_0\in G\setminus\{O,c_*P\}$,
four points $U_n,U_s,U_t,U_w\in G$, and one rational slope $\eta\ne0$.
The same point $H_0$ is used in both parts below.

\paragraph{The scalar condition.}
Set
\begin{equation}\label{sd:scalar}
 \begin{split}
 W_\beta={}&\left(1+\frac{\rho(Q_s)^2}{R_0}\right)
             \left(1+\frac{\rho(H_0)^2}{R_0}\right)
             \left(1+\frac{(z(Q_s)/z(H_0))^2}{R_0^{2c-1}}\right),\\
 \Theta_\beta={}&\sdTr(W_\beta^{-1}),\qquad
 W_\beta\Theta_\beta\ne0,\qquad
 \Phi_0\left(\frac{R_0}{(N\Theta_\beta)^2}\right).
 \end{split}
\end{equation}
This is the aggregation \eqref{sd:agg} for $R=R_0/N^2$ and the three
expressions $D_j/N^{2k_j-1}$, where
$(D_j,k_j)=(\rho(Q_s),1),(\rho(H_0),1),(z(Q_s)/z(H_0),c)$.
The normalizing powers cancel from the product, leaving the shorter
displayed formula.

\paragraph{The main conditions.}
For $i\in\{n,s,t,w\}$ set
\[
 V_0=H_0-c_*P,\quad V_i=U_i-[c_*]Q_i,\qquad
 L_i=\frac{\ell_K(z(U_i))}{\ell_K(z(H_0))},\quad
 L_i'=\frac{\ell_K^\sigma(z(V_i^\sigma))}
                    {\ell_K^\sigma(z(V_0^\sigma))}.
\]
Both logarithm denominators are required to be nonzero. Form the following
sixteen expressions with their precisions:
\begin{equation}\label{sd:list}
 \begin{array}{c|c|c}
 \text{expressions}&\text{number}&\text{precision}\\\hline
 \rho(H_0),\ \rho(U_i)&5&1\\
 \rho(V_0^\sigma),\ \rho(V_i^\sigma)&5&1\\
 L_i-L_i'&4&K\\
 g_L(L_s,L_t,L_w),\ \beta-L_n&2&K.
 \end{array}
\end{equation}
Denote this list by $(D_j,k_j)_{j=1}^{16}$.
Write the eight rational components of the four $x(Q_i)$ as $X_1,\ldots,X_8$.
Set $a_0=1$, and $a_r=X_r^2$ when $X_r\ne0$, $a_r=1$ otherwise.
For $0\le r\le8$ put
\[
 t_r=a_r\eta,\qquad
 R_r=a_r^{-5}\frac{\prod_{j=0}^4(t_r-\theta_j)}{(t_r-\xi)^6}.
\]
Require $t_r\notin\{\theta_0,\ldots,\theta_4,\xi\}$ and
\begin{equation}\label{sd:main}
 \mathcal A\left(R_r;
       \left(\frac{D_j}{(t_r-\xi)^{k_j+1}},k_j\right)_{j=1}^{16}\right)
 \qquad(0\le r\le8).
\end{equation}
All equalities over $F$, rational-function evaluations, finite branches and
nonzero guards are interpreted on their usual finite rational charts.
Equations \eqref{sd:phi}, \eqref{sd:agg}, \eqref{sd:scalar}, and
\eqref{sd:main}, together with the point and denominator conditions, are a
finite existential formula over $\sdQ$. No point index is quantified as a
rational number.

\subsection{Soundness}

\begin{sdlemma}\label{sd:scalar-height}
If \eqref{sd:scalar} holds and $Q_s=sP$, then $h(\beta)\le H|s|$.
\end{sdlemma}
\begin{proof}
Let $q\mid M_S(\beta)$. At a finite normalized contact of positive depth,
$\beta$ is $q$-integral and exactly one $\beta-\theta_j$ has positive odd
valuation. Thus $v_q(r)>0$ is odd, $N$ is a $q$-unit, and
$v_q(R_0/N^2)=v_q(r)$ is odd. By \eqref{sd:agg-sound}, at some place over
$q$ both points are formal and $v(z(Q_s)/z(H_0))\ge c$.
Writing $H_0=hP$, \eqref{sd:formal} gives $v_q(s/h)\ge c$ and hence $q^c\mid s$.
Therefore $M_S(\beta)^c\mid s$, and \eqref{sd:height} proves the result.
If $v_q(\beta)<0$, all five normalized finite-contact depths are zero:
the numerator of each primitive form is a unit. Such a prime does not
occur in $M_S(\beta)$. Thus no inference involving the normalization $N$
at a negative valuation is used.
\end{proof}

\begin{sdlemma}\label{sd:main-errors}
Write $Q_i=n_iP$ and $B=\max_i|n_i|$. If the main conditions hold, then
\[
 B^2\le C|D|^{1/2}
\]
whenever either $D=g_L(n_s,n_t,n_w)\ne0$, or
$D=\sdnum(\beta-n_n)\ne0$.
\end{sdlemma}
\begin{proof}
Fix $q\mid M_S(t_r)$. Then $t_r$ is integral, exactly one finite contact
has positive odd depth, and $t_r-\xi$ is a unit. Since $a_r^{-5}$ is a
rational square, $v_q(R_r)$ is odd. Aggregation supplies one place at which
the original sixteen expressions in \eqref{sd:list} have their specified
precisions.

Write $H_0=hP$, $U_i=z_iP$. The first ten tests make the relevant points
formal, and the next four give $L_i-L_i'\in q^K\mathcal O_v$.
The integers $h,h-c_*$ are nonzero and at least one is a $q$-unit.
Consequently at least one of $z_i/h,(z_i-c_*n_i)/(h-c_*)$ is integral.
The valuation identity for truncated logarithms and $L_i-L_i'$ then make
both ratios integral. The congruence in \eqref{sd:formal} applies to both.
Multiplying by $h,h-c_*$ and subtracting yields
\[
 c_*n_i\equiv hL_i-(h-c_*)L_i'\equiv c_*L_i\pmod{q^K\mathcal O_v}.
\]
Thus $L_i\equiv n_i\pmod{q^K\mathcal O_v}$. The final two tests show
that $q^K$ divides either nonzero integer $D$ under consideration. For the
second test, $\beta-L_n$ first makes $\beta$ integral at $q$, so the reduced
denominator of $\beta-n_n$ is a $q$-unit.

Hence $M_S(t_r)^K\mid D$ for every $r$, and $K=2c$ gives
$h(t_r)\le H|D|^{1/2}$. Since
$2h(X_r)=h(a_r)\le h(\eta)+h(a_r\eta)$ for nonzero $X_r$,
every component satisfies $h(X_r)\le H|D|^{1/2}$; a zero component has
height zero. For a point whose index has absolute value $B$,
\[
 2B^2\widehat h(P)
 \le h(x(Q_i))+C_E
 \le 2H|D|^{1/2}+\tfrac32\log2+C_E.
\]
Use $|D|\ge1$ and \eqref{sd:C}.
\end{proof}

\begin{sdproposition}\label{sd:soundness}
The displayed formula implies $X\in\sdZ$.
\end{sdproposition}
\begin{proof}
Write the four indices as $n,s,t,w$, and put $B=\max\{|n|,|s|,|t|,|w|\}$.
They are all nonzero. The integer residual satisfies
$|g_L(s,t,w)|\le C_LB^4$. If nonzero, Lemma~\ref{sd:main-errors} and
\eqref{sd:C} give $B^2<B^2/2$. Hence $g_L(s,t,w)=0$.

Jones's Lucas classification \cite[Corollary~2.7]{JonesLucas2003} gives
\[
 |t|\ge(|Ls|-1)^{|Ls|-4}.
\]
For completeness, write $a=|Ls|\ge9$ and
$y=\operatorname{sgn}(Ls)(Ls-2)w$. The equation becomes
$(2t-ay)^2-(a^2-4)y^2=4$. Jones's classification and the recurrence
$U_0=0,U_1=1,U_{j+1}=aU_j-U_{j-1}$ give
$|y|=U_k$, $|t|\in\{U_{k-1},U_{k+1}\}$.
The congruences $U_k\equiv k\pmod{a-2}$ and
$U_k\equiv(-1)^{k-1}k\pmod{a+2}$ imply $k\ge a-2$;
$U_j\ge(a-1)^{j-1}$ gives the assertion. Since $L>H$ and $L\ge9$,
\begin{equation}\label{sd:growth}
 \log B\ge(L|s|-4)\log(L|s|-1)>H|s|.
\end{equation}

If $\beta\ne n$, let $D=\sdnum(\beta-n)$. Then
$|D|\le2B\exp h(\beta)$. Lemmas~\ref{sd:scalar-height} and
\ref{sd:main-errors} imply
\[
 3\log B\le 2\log C+\log2+h(\beta)<H|s|,
\]
contrary to \eqref{sd:growth}. Thus $\beta=n\in\sdZ$, so
$X^2=\beta-b_*\in\sdZ$ and $X\in\sdZ$.
\end{proof}

\subsection{Completeness and sharing the local anchor}

We give the order of choices explicitly. Fix $X\in\sdZ$.
Then $\beta\in\sdZ$, and $r=r(\beta)$ is a positive integer.
Let $\mathcal P$ be the set of primes dividing $Ar$. For $p\in\mathcal P$
put $\alpha_p=2v_p(A)+v_p(r)>0$. The integer $N=1+T_0^2r^2$ is a unit
at every such prime: this is immediate at divisors of $r$, and at divisors
of $A$ it follows from the divisibility imposed on $T_0$.

Choose a positive integer $\lambda$ such that
\[
 v_p(\lambda)\ge c\alpha_p+v_p(2)+1\qquad(p\in\mathcal P).
\]
Choose a positive integer $D_0$, divisible by $c_*$, so that every point
in $D_0G$ belongs, at all places over $\mathcal P$, to a sufficiently deep
formal subgroup and satisfies
$v(\rho(Q))\ge\alpha_p+v_p(2)+1$ unless $Q=O$.
Such a $D_0$ exists by the finite index of deep local formal subgroups.
Choose a positive multiple $j_0$ of $D_0$ for which
$Q_s=[\lambda j_0]P$ belongs to one of the open sets
\eqref{sd:goodsets} at the real places. This is possible with arbitrarily
large $j_0$: the closure of any nonconstant arithmetic progression in
$\mathcal G$ contains an entire component, because multiplication by a
nonzero integer is surjective on its identity torus. The positive tail of
the progression has the same closure.

Put $s=\lambda j_0$, $n=\beta$, and choose positive $t,w$ satisfying
$g_L(s,t,w)=0$. Explicitly take the Lucas sequence with parameter $Ls$,
put $k=Ls-2$, and take $t=U_{k+1}$ and $w=U_k/k$.
The congruence $U_k\equiv k\pmod{Ls-2}$ and the identity
$U_{k+1}^2-LsU_{k+1}U_k+U_k^2=1$ verify these choices.
All four root points are now fixed.

Set
\begin{equation}\label{sd:fixedmodulus}
 D_1=\operatorname{lcm}\left(D_0,
               \prod_{p\in\mathcal P}p^{v_p(j_0)+1}\right).
\end{equation}
We will require the shared index $h$ to satisfy $h\equiv j_0\pmod{D_1}$.
Then $H_0=hP$ has the same formal depths as $j_0P$ at the places in
$\mathcal P$, and
\begin{equation}\label{sd:scalar-depths}
 v(\rho(H_0)),v(\rho(Q_s))\ge\alpha_p+v_p(2)+1,
 \quad v(z(Q_s)/z(H_0))=v_p(\lambda)\ge c\alpha_p+v_p(2)+1.
\end{equation}
Indeed both $hP,j_0P$ lie in the chosen deep formal subgroup and
$v_p(h)=v_p(j_0)$, so the formal logarithm compares their parameter
valuations. Also $c_*\mid h$.

We additionally choose $H_0$ in an open set from \eqref{sd:goodsets}.
Then the three unnormalized scalar expressions have absolute values at
most $C_1$ at both embeddings. At these embeddings the three nonconstant
terms in the scalar product \eqref{sd:agg} have the exact form
\[
 D_j^2 A^{2-4k_j}r^{1-2k_j},
 \qquad (D_j,k_j)=(\rho(Q_s),1),(\rho(H_0),1),(z(Q_s)/z(H_0),c).
\]
They are nonnegative and at most $C_1^2$.
Consequently $1\le W_\beta\le(1+C_1^2)^3$,
$\Theta_\beta\ge2/(1+C_1^2)^3$, and
\[
 0<\frac{R_0}{(N\Theta_\beta)^2}
 \le\frac{A^2(1+C_1^2)^6}{4T_0^4}<\epsilon.
\]
At $p\in\mathcal P$, \eqref{sd:scalar-depths} gives valuation at least
$\alpha_p+2v_p(2)+2$ for each nonconstant term. Thus
$\Theta_\beta=2+\text{an element of valuation }>v_p(2)$ and
\[
 v_p\left(R_0/(N\Theta_\beta)^2\right)=\alpha_p-2v_p(2).
\]
This is at least $b_p$ on $S_0$ by \eqref{sd:A}, and is nonnegative at
every other prime of $\mathcal P$. Outside $\mathcal P$ the valuation of
$R_0/N^2$ is even, so that of $R_0/(N\Theta_\beta)^2$ is even.
The input therefore has no negative odd valuation at any prime.
The completeness assertion for $\Phi_0$ proves the scalar condition,
for every later $h$ having the stipulated congruence and real bounds.

It remains to construct such an $h$ together with the main slope.
Set $U_i=[n_i]H_0$ for $i=n,s,t,w$ in the completeness construction.
Then $V_i=[n_i]V_0$. The basic reciprocal coordinates and each individual
truncated-logarithm ratio are rational functions on one real elliptic factor.
Their poles and logarithm-denominator zeros lie over finite subsets of
that factor. Bounds for these basic functions bound all sixteen expressions
in \eqref{sd:list}. None of the forbidden finite fibres contains a
component of $\mathcal G$.
Inside the previously chosen open set in each component, remove these
fibres and choose a smaller nonempty open set. On these smaller sets all
sixteen main expressions have a common bound $C_2$. Crucially $C_2$ is
chosen before the main prime pattern and is independent of its later
congruence moduli.

At every place over $S_0$, the sixteen main expressions are integral
whenever $c_*\mid h$ and $h\notin\{0,c_*\}$. Indeed $hP,(h-c_*)P$ lie
in the fixed deep formal subgroup. Their integer multiples give truncated
logarithm ratios of valuation $v_p(n_i)\ge0$, and their reciprocal
$x$-coordinates are integral. Since $\beta$ and the coefficients of
$g_L$ are integers, the last two expressions are integral as well.

For the now fixed positive squares $a_r$, choose a square denominator
$v=d_2^2w^2$ with $d_2$ sufficiently divisible and put $\eta=u/v$.
Clear denominators so that all distinct positions
\begin{equation}\label{sd:positions}
             u-(\theta_j/a_r)v
       \quad(0\le r\le8,\ 0\le j\le4)
\end{equation}
are integral quadratic polynomials in $u,w$. Since $\theta_0=0$, there
are at most $37$ distinct positions. Apply the prime-pattern input with
the odd squarefree product of the primes dividing $D_1$, except $7$, as
its avoidance modulus. Take all selected primes beyond the fixed
exceptions and the supports and differences of these rational coefficients.
The theorem also permits arbitrarily large $\eta$ and arbitrary fixed
denominator divisibility.

Let $U$ be the least common multiple of all selected $d_+(q)$ and $V$ that
of all selected $d_-(q)$. Impose
\begin{equation}\label{sd:CRT}
 h\equiv0\pmod U,\qquad h\equiv c_*\pmod V,
 \qquad h\equiv j_0\pmod{D_1}.
\end{equation}
These congruences are compatible. Indeed $\gcd(U,V)=4\mid c_*$.
Every common divisor of $UV$ and $D_1$ is a power of two: the avoidance
condition handles all odd primes other than $7$, and the fixed Gaussian
residue handles $7$. The relevant powers of two divide $4$ for $U$ and
$64$ for $V$. Both $j_0$ and $c_*$ are divisible by $64$.
This checks every pair of congruences, including arbitrary powers in
$D_1$. The generalized Chinese remainder theorem therefore gives a
nonconstant arithmetic progression of solutions. Choose an arbitrarily
large positive solution whose real point belongs to one of the smaller
open sets defining $C_2$; avoid $h=c_*$.

At each selected contact prime orient the two places so that the first
has order $d_+(q)$. Then $H_0,U_i$ are formal at the first place and
$V_0,V_i$ at the second. The two rational index ratios in the formal
comparison are both exactly $n_i$. Thus all sixteen main expressions
have the required precisions at the first place: conjugation supplies
the second family there, the logarithm differences vanish modulo $q^K$,
and $g_L(n_s,n_t,n_w)=\beta-n_n=0$.

For completeness of the pole predicates, the exact detector formula is
\[
 R_r=\frac{\prod_{j=0}^4\bigl(u-(\theta_j/a_r)v\bigr)}
                    {v^5(t_r-\xi)^6}.
\]
Within each slope its five contact primes are distinct; $v$ is a square
and is a unit at those primes. Hence $v_q(R_r)=1$ at its five contacts,
and $v_p(R_r)$ is even elsewhere. At a contact, every local factor in
$W$ is either $1$ modulo its place or has negative valuation. At the
successful place they are all $1$ modulo that place. Consequently the
trace of $W^{-1}$ has residue $1$ or $2$, which is nonzero, so the pole
input has valuation one there and even valuation elsewhere. It has no
negative odd valuation.

Here are the bounds ensuring the remaining preparation requirements.
Uniformly when the sixteen unscaled expressions are bounded by $C_2$,
as $\eta\to+\infty$,
\[
 \begin{gathered}
 R_r=a_r^{-6}\eta^{-1}(1+O(\eta^{-1})),\quad
 W_r=1+O(\eta^{-3}),\quad \Theta_r=2+O(\eta^{-3}),\\
 R_r/\Theta_r^2=(4a_r^6\eta)^{-1}(1+O(\eta^{-1})).
 \end{gathered}
\]
Thus all nine pole inputs are positive and less than $\epsilon$ once a
fixed lower bound on $\eta$ is imposed before the pattern is selected.
For $p\in S_0$ put $\nu=v_p(v)$ and $\gamma_r=v_p(a_r)$, both even.
For sufficiently divisible $v$, the prime $u$ is a $p$-unit and
\[
 v_p(R_r)=\nu-6\gamma_r,\qquad
 v\left(R_r\left(\frac{D_j/(t_r-\xi)^{k_j+1}}{R_r^{k_j}}\right)^2\right)
 \ge3\nu+(10k_j-8)\gamma_r.
\]
All unscaled $D_j$ are integral. Increasing the prescribed divisibility
of $v$ makes the right side exceed $v_p(2)$ for every $r,j$.
It follows that $v_p(\Theta_r)=v_p(2)$ and
$v_p(R_r/\Theta_r^2)=\nu-6\gamma_r-2v_p(2)\ge b_p$.
All bounds were fixed before the prime pattern and its CRT modulus.
The completeness assertion for $\Phi_0$ now verifies every main predicate.
This completes the proof that the formula holds for every integer $X$.

\subsection{Polynomial compilation and variable count}

\begin{sdtheorem}\label{sd:forty}
Assume the arithmetic inputs above. There is a polynomial
$F\in\sdZ[X,Y_1,\ldots,Y_{40}]$ such that, for every $X\in\sdQ$,
\[
 X\in\sdZ\quad\Longleftrightarrow\quad
 \exists Y_1,\ldots,Y_{40}\in\sdQ\quad F(X,\mathbf Y)=0.
\]
Given numerical realizations of the fixed constants and finite multiplier
data above, the proof gives a terminating finite algorithm to produce $F$.
It does not claim that this polynomial, its circuit, or its coefficients
have been emitted here.
\end{sdtheorem}
\begin{proof}
The equivalence was proved above. There are nine cyclic points: four root
points and five shared local points. Represent each as $[m](x,y)$ with
$(x,y)\in E(F)$ affine. Every cyclic point admits such an affine preimage;
for $O$, use the two-torsion point $(0,0)$.
Writing $x=u+v\sqrt2$ gives two free rational coordinates. If
$x^3-d^2x=C_0+C_1\sqrt2$ and $C_1\ne0$, its ordinate is encoded by the
single finite choice
\[
 \tau^4-C_0\tau^2+C_1^2/2=0,\qquad
 y=\tau+\frac{C_1}{2\tau}\sqrt2.
\]
This quartic is separable and $\tau\ne0$. If $C_1=0$, use the finite
branches $y=\tau$, $\tau^2=C_0$, or $y=\tau\sqrt2$, $\tau^2=C_0/2$;
split off $C_0=0$. Thus each point contributes two free coordinates and
one finite algebraic choice.

There are ten pole predicates: nine main predicates and one scalar
predicate. Each contributes two structural coordinates $e,z$ and two
finite square-root choices. There is one free slope $\eta$.
Consequently the free structural parameters number
\[
                     9\cdot2+10\cdot2+1=39.
\]
The finite-algebra compilation of \cite[Section~10]{Sun2026} combines
every remaining finite choice, branch and nonzero guard into one final
rational variable. For clarity, it first separates the points of each
finite reduced algebra by one of finitely many fixed linear combinations,
retains a nonzero discriminant, recovers the component roots by matrix
inversion, and filters equalities and inequations by polynomial gcds.
Dependent square-root choices are then incorporated in a second reduced
algebra and separated again. Degree drops and zero square inputs are
handled by finite branches. Finally the reciprocal-polynomial construction
absorbs every denominator and nonzero guard into that same variable;
sums of squares combine equations and products combine the finitely many
branches. This is an algorithm on polynomials of fixed finite degree and
uses no additional free parameters. It gives $39+1=40$ existential variables.
\end{proof}

There is a less compressed presentation that avoids the global finite-algebra
step. Keep one ordinate variable for each point. For the two square tests
in a pole predicate, use the elementary identity
\[
 A,B\in\sdSq\quad\Longleftrightarrow\quad
 \exists t\in\sdQ:\quad
 (A^2+B^2)\bigl(1-2(A+B)t^2+(A-B)^2t^4\bigr)=0.
\]
Indeed, when not both inputs vanish, a rational root gives a nonzero
rational number equal to a signed sum of square roots, from which both
square roots are recovered using their sum and difference; the converse
chooses a nonzero signed sum. This leaves
\[
              9\cdot3+10\cdot3+1+1=59
\]
variables, the last enforcing the product of all nonzero guards on each
finite branch. The formula consists of the displayed rational-function
compositions, local quartics and finite branches. It is likely more useful
for a short circuit description than the forty-variable polynomial.
Before branching, the uncompressed finite fibres already have degree at
most $4^9 2^{20}=2^{38}$, so global primitive-element packing can increase
degree and coefficient data enormously. Neither count alone establishes
an improvement in total encoded bit length.
